\documentclass{beamer}
\usetheme{Madrid}

\title{GENMANIP: LLM-driven Simulation for Generalizable}
\author{Ning Gao1}
\date{}

\begin{document}

\frame{\titlepage}

\begin{frame}{GENMANIP: LLM-driven Simulation for Generalizable}
Authors: Ning Gao1, 2∗, Yilun Chen1∗‡, Shuai Yang1, 3∗, Xinyi Chen1, 4∗, Yang Tian1, Hao Li1
\end{frame}

\begin{frame}{Abstract}
Robotic manipulation in real-world settings remains chal-
lenging, especially regarding robust generalization. Exist-
ing simulation platforms lack sufficient support for explor-
ing how policies adapt to varied instructions and scenarios.
Thus, they lag behind the growing interest in instruction-
following foundation models like LLMs, whose adaptability
is crucial yet remains underexplored in fair comparisons.
To bridge this gap, we introduce GENMANIP, a realistic
tabletop simulation platform t
\end{frame}

\begin{frame}{Abstract}
Robotic manipulation in real-world settings remains chal-
lenging, especially regarding robust generalization. Exist-
ing simulation platforms lack sufficient support for explor-
ing how policies adapt to varied instructions and scenarios.
Thus, they lag behind the growing interest in instruction-
following foundation models like LLMs, whose adaptability
is crucial yet remains underexplored in fair comparisons.
To bridge this gap, we introduce GENMANIP, a realistic
tabletop simulation platform tailored for policy generaliza-
tion studies. It features an automatic pipeline via LLM-
driven task-oriented scene graph to synthesize large-scale,
diverse tasks using 10K annotated 3D object assets. To sys-
tematically assess generalization, we present GENMANIP-
BENCH, a benchmark of 200 scenarios refined via human-
in-the-loop corrections. We evaluate two policy types: (1)
modular manipulation systems integrating foundation mod-
els for perception, reasoning, and planning, and (2) end-
to-end 
\end{frame}

\begin{frame}{Introduction}
Policy generalization is a core challenge within the field of
robotic manipulation. Scholars have identified two popular
approaches to enhance the generalizability of models. One
streams [29, 30, 32, 44, 51] involves the hierarchical use of
established foundation models, which utilize multi-modal
large language models [3, 7, 23, 36, 54, 59] (MLLMs), and
∗Equal contribution.
‡ Project lead.
† Corresponding author.
2D foundation models [36, 49, 57, 59] for perception, rea-
soning, or task and moti
\end{frame}

\begin{frame}{Introduction}
Policy generalization is a core challenge within the field of
robotic manipulation. Scholars have identified two popular
approaches to enhance the generalizability of models. One
streams [29, 30, 32, 44, 51] involves the hierarchical use of
established foundation models, which utilize multi-modal
large language models [3, 7, 23, 36, 54, 59] (MLLMs), and
∗Equal contribution.
‡ Project lead.
† Corresponding author.
2D foundation models [36, 49, 57, 59] for perception, rea-
soning, or task and motion planning. Another strategies
emphasizes end-to-end imitation learning, including recent
popular vision-language-action models (VLAs) pretrained
with large-scale vision-language models (VLMs). trained
on large-scale robotic datasets [53] and accomplish down-
stream tasks in a zero-shot or downstream fine-tuning man-
ner. However, benchmarking both policy types for general-
ization in real-world environments remains challenging due
to the high cost and difficulty of replicating exact scenarios

\end{frame}

\begin{frame}{Related Work}
s
Simulation platforms for robotics. Significant advance-
ments have been made in developing robot simulators [33,
37, 46]. Many simulation frameworks [19, 22, 35, 48, 55,
61, 78, 80] have been built for robotics based on these sim-
ulators. But very few of them can stratify the requirement
of diversity in the type of scenes, objects and activities, and
realism of the underlying simulation environments as men-
tioned by Behavior-1K [39]. Moreover, generating demon-
strations is crucial in the fi
\end{frame}

\begin{frame}{Related Work}
s
Simulation platforms for robotics. Significant advance-
ments have been made in developing robot simulators [33,
37, 46]. Many simulation frameworks [19, 22, 35, 48, 55,
61, 78, 80] have been built for robotics based on these sim-
ulators. But very few of them can stratify the requirement
of diversity in the type of scenes, objects and activities, and
realism of the underlying simulation environments as men-
tioned by Behavior-1K [39]. Moreover, generating demon-
strations is crucial in the field of robot learning.
Many
methods have been developed to generate demonstration,
including rule-based finite state machine [28] and transfer-
ing human demonstration [47, 50]. Differently, GENMA-
NIP-BENCH incorporates hand-designed skills such as mo-
tion planner and grasping [14] by leveraging privileged in-
formation to generate feasible demonstrations.
Task and scenario generation.
Scenario generation is
widely used in robotics [8, 39, 50, 69] and autonomous driv-
ing [10, 11, 17] to creat
\end{frame}

\begin{frame}{Experiments}
Our experiments aim to address two key questions: (1)
What insights can a well-designed benchmark provide
about current modular VLM agents? (2) Can learning-based
methods generalize across varying task complexities and
benefit from data scaling?
5.1. Main Results
We evaluate two prompt-based methods, MOKA [44] and
CoPA [30], within our benchmark framework, with detailed
results in Tab. 2. CoPA, leveraging GPT-4.5’s capabilities,
achieves a SR of 23.0% across all tasks, significantly out-
perform
\end{frame}

\begin{frame}{Experiments}
Our experiments aim to address two key questions: (1)
What insights can a well-designed benchmark provide
about current modular VLM agents? (2) Can learning-based
methods generalize across varying task complexities and
benefit from data scaling?
5.1. Main Results
We evaluate two prompt-based methods, MOKA [44] and
CoPA [30], within our benchmark framework, with detailed
results in Tab. 2. CoPA, leveraging GPT-4.5’s capabilities,
achieves a SR of 23.0% across all tasks, significantly out-
performing other models. This demonstrates CoPA’s supe-
rior performance in handling diverse tasks. However, the
evaluation also highlights the challenges posed by long-
horizon tasks, which emerge as the most difficult category.
All models, including CoPA and MOKA, achieve an aver-
age success rate of 9.07% on these tasks. This underscores
a limitation of prompt-based models in managing the split-
ting of complex tasks and planning across a long horizon.
In contrast, for tasks that require understandi
\end{frame}

\begin{frame}{Conclusion}
We introduced GENMANIP, a large-scale simulation frame-
work for benchmarking generalist robotic manipulation
across diverse scenes and instructions.
Central to it is
the task-oriented scene graph (ToSG), a scalable, natu-
ral language–based representation that models object-level
and spatial relations for flexible task generation. We also
present GENMANIP-BENCH, a set of 200 realistic scenar-
ios for evaluating generalization. Experiments show that
modular methods generalize better but struggle
\end{frame}

\begin{frame}{Conclusion}
We introduced GENMANIP, a large-scale simulation frame-
work for benchmarking generalist robotic manipulation
across diverse scenes and instructions.
Central to it is
the task-oriented scene graph (ToSG), a scalable, natu-
ral language–based representation that models object-level
and spatial relations for flexible task generation. We also
present GENMANIP-BENCH, a set of 200 realistic scenar-
ios for evaluating generalization. Experiments show that
modular methods generalize better but struggle with spatial
perception, while end-to-end policies benefit from scaling
but remain sensitive to scene and instruction changes.
\end{frame}

\begin{frame}{References}
[1] Constructions Aeronautiques, Adele Howe, Craig Knoblock,
ISI Drew McDermott, Ashwin Ram, Manuela Veloso, Daniel
Weld, David Wilkins Sri, Anthony Barrett, Dave Christian-
son, et al. Pddl— the planning domain definition language.
Technical Report, Tech. Rep., 1998. 1, 4
[2] Peter Anderson, Angel Chang, Devendra Singh Chaplot,
Alexey Dosovitskiy, Saurabh Gupta, Vladlen Koltun, Jana
Kosecka, Jitendra Malik, Roozbeh Mottaghi, Manolis Savva,
et al. On evaluation of embodied navigation agents. arX
\end{frame}

\begin{frame}{References}
[1] Constructions Aeronautiques, Adele Howe, Craig Knoblock,
ISI Drew McDermott, Ashwin Ram, Manuela Veloso, Daniel
Weld, David Wilkins Sri, Anthony Barrett, Dave Christian-
son, et al. Pddl— the planning domain definition language.
Technical Report, Tech. Rep., 1998. 1, 4
[2] Peter Anderson, Angel Chang, Devendra Singh Chaplot,
Alexey Dosovitskiy, Saurabh Gupta, Vladlen Koltun, Jana
Kosecka, Jitendra Malik, Roozbeh Mottaghi, Manolis Savva,
et al. On evaluation of embodied navigation agents. arXiv
preprint arXiv:1807.06757, 2018. 5
[3] Anthropic. Claude ai, 2025. 1
[4] Homanga Bharadhwaj, Roozbeh Mottaghi, Abhinav Gupta,
and Shubham Tulsiani. Track2act: Predicting point tracks
from internet videos enables diverse zero-shot robot manip-
ulation, 2024. 3
[5] Anthony Brohan, Noah Brown, Justice Carbajal, Yevgen
Chebotar, Joseph Dabis, Chelsea Finn, Keerthana Gopalakr-
ishnan, Karol Hausman, Alex Herzog, Jasmine Hsu, et al.
Rt-1: Robotics transformer for real-world control at scale.
arXiv 
\end{frame}

\begin{frame}{Figures - Figures 1-2}
\begin{figure}
  \centering
  \includegraphics[width=0.8\textwidth]{figure_1.png}
\end{figure}
\begin{figure}
  \centering
  \includegraphics[width=0.8\textwidth]{figure_2.png}
\end{figure}
\end{frame}

\begin{frame}{Figures - Figures 3-4}
\begin{figure}
  \centering
  \includegraphics[width=0.8\textwidth]{figure_3.png}
\end{figure}
\begin{figure}
  \centering
  \includegraphics[width=0.8\textwidth]{figure_4.png}
\end{figure}
\end{frame}

\begin{frame}{Figures - Figures 5-6}
\begin{figure}
  \centering
  \includegraphics[width=0.8\textwidth]{figure_5.png}
\end{figure}
\begin{figure}
  \centering
  \includegraphics[width=0.8\textwidth]{figure_6.png}
\end{figure}
\end{frame}

\begin{frame}{Tables - Tables 1-2}
\begin{table}
  \centering
  \includegraphics[width=0.9\textwidth]{table_1.png}
\end{table}
\begin{table}
  \centering
  \includegraphics[width=0.9\textwidth]{table_2.png}
\end{table}
\end{frame}

\begin{frame}{Tables - Tables 3-4}
\begin{table}
  \centering
  \includegraphics[width=0.9\textwidth]{table_3.png}
\end{table}
\begin{table}
  \centering
  \includegraphics[width=0.9\textwidth]{table_4.png}
\end{table}
\end{frame}

\end{document}
